%
% teil3.tex -- Beispiel-File für Teil 3
%
% (c) 2020 Prof Dr Andreas Müller, Hochschule Rapperswil
%
% !TEX root = ../../buch.tex
% !TEX encoding = UTF-8
%
\section{Die Gamma-Funktion und ihre Produktdarstellung}
\kopfrechts{Die Gamma-Funktion und ihre Produktdarstellung}

\subsection{Definition und erste Eigenschaften}

Dieser Abschnitt ist wiederum eine vereinfachte Betrachtung, 
für eine detaillierte Analyse siehe Kapitel \ref{chapter:gamma}.
Die Gamma-Funktion ist die natürliche Verallgemeinerung der Fakultät
\index{Gamma-Funktion}%
auf komplexe Argumente.
Für $\operatorname{Re}(z) > 0$ ist sie durch das Integral
\begin{equation}
\Gamma(z)
=
\int_0^{\infty} t^{z-1} e^{-t} \, dt
\end{equation}
\index{G(z)@$\Gamma(z)$}%
definiert.
Es gilt $\Gamma(n+1) = n!$ für natürliche Zahlen $n$,
sowie die Funktionalgleichung $\Gamma(z+1) = z\,\Gamma(z)$,
die eine meromorphe Fortsetzung auf ganz $\mathbb{C}$ erlaubt.
Die Gamma-Funktion hat keine Nullstellen, besitzt aber einfache Pole
bei $z = 0, -1, -2, \ldots$

Dieser Umstand ist entscheidend:
Man schreibt die Produktdarstellung nicht für $\Gamma(z)$ selbst,
sondern für $1/\Gamma(z)$, die eine ganze Funktion ist.
Dies ist ein häufiger Fallstrick beim ersten Studium der Gamma-Funktion.

\subsection{Die Weierstraß-Produktdarstellung}

Da $1/\Gamma(z)$ eine ganze Funktion mit einfachen Nullstellen bei
$z = 0, -1, -2, \ldots$ ist, lässt sich der Weierstraß-Produktsatz
direkt anwenden.
Mit den Elementarfaktoren $E_1(z) = (1-z)e^z$ ergibt sich die Formel
\cite{produkt:steinShakarchi}:
\begin{equation}
\frac{1}{\Gamma(z)}
=
z \cdot e^{\gamma z} \cdot
\prod_{n=1}^{\infty} \biggl(1 + \frac{z}{n}\biggr) e^{-z/n},
\end{equation}
wobei $\gamma \approx 0{,}5772$ die Euler-Mascheroni-Konstante ist.
\index{Euler-Mascheroni-Konstante}%
Diese Konstante tritt als der Koeffizient in $g(z) = \gamma z$ auf,
der aus dem Hadamard-Satz und der Normierungsbedingung $\Gamma(1) = 1$ bestimmt wird.

Die Produktdarstellung macht die analytische Struktur der Gamma-Funktion
vollständig transparent:
Die Pole bei den nicht-positiven ganzen Zahlen sind direkt als Nullstellen
von $1/\Gamma(z)$ ablesbar, und die Euler-Mascheroni-Konstante erscheint
als natürliche Normierungsgröße.

\subsection{Die Reflexionsformel}

Die eleganteste Verbindung zwischen der Gamma-Funktion und dem Sinus
liefert die \emph{Reflexionsformel von Euler}:
\index{Reflexionsformel!für $\Gamma(z)$}%
\begin{equation}
\Gamma(z) \cdot \Gamma(1-z)
=
\frac{\pi}{\sin(\pi z)}.
\end{equation}
Diese Formel lässt sich mit den Produktdarstellungen beider Seiten rigoros beweisen.
Der Beweis lässt sich in drei Schritte gliedern:
Erstens betrachtet man $f(z) = \Gamma(z) \cdot \Gamma(1-z) \cdot \sin(\pi z)$
und zeigt, dass $f$ nach Multiplikation eine ganze Funktion ist.
Zweitens ergibt die Wachstumsanalyse (mithilfe der Stirlingschen Formel
oder der Produktdarstellung), dass $f$ beschränkt und damit nach dem
Satz von Liouville konstant ist.
Drittens bestimmt die Normierungsbedingung bei $z = 1/2$ die Konstante zu $\pi$,
woraus die Formel folgt.

Die Reflexionsformel ist mehr als eine Kuriosität.
Sie zeigt, dass die Produktdarstellung des Sinus und die Produktdarstellung
von $1/\Gamma(z)$ zwei Seiten derselben Münze sind,
denn beide sind Manifestationen des Weierstraß-Produktsatzes auf denselben
Gebieten der komplexen Ebene.
Als spezieller Fall liefert $z = 1/2$ die schöne Identität:
\begin{equation}
\Gamma\!\biggl(\frac{1}{2}\biggr) = \sqrt{\pi}.
\end{equation}

\section{Fazit und Ausblick}
\kopfrechts{Fazit und Ausblick}

\subsection{Mathematische Einordnung}

Der Weierstraß-Produktsatz zeigt, dass die Nullstellenstruktur eine ganze Funktion
fast vollständig bestimmt, der verbleibende Freiheitsgrad ist allein der
nullstellenfreie Faktor $e^{g(z)}$.
Das ist das präzise analytische Analogon zur algebraischen Aussage,
dass ein Polynom durch seine Nullstellen und seinen Leitkoeffizienten
eindeutig festgelegt ist.

Für ganze Funktionen endlicher Ordnung schärft der Hadamard-Faktorisierungssatz
dieses Bild weiter:
$g(z)$ ist dann ein Polynom vom Grad höchstens $\lfloor \varrho \rfloor$.
In den beiden wichtigsten Beispielen dieser Arbeit
$\sin(\pi z)$ und $1/\Gamma(z)$ ist $\varrho = 1$,
sodass $g$ linear ist und durch Normierungsbedingungen vollständig bestimmt wird.

\subsection{Praktische Anwendungen}

Die hier entwickelte Theorie mag abstrakt wirken, hat aber weitreichende
Konsequenzen in Wissenschaft und Technik.

\subsubsection{Signalverarbeitung und Regelungstechnik}
Ein lineares Filter wird vollständig durch seine Übertragungsfunktion beschrieben,
die ihrerseits eine rationale Funktion ist und durch ihre Pole und Nullstellen
(bis auf einen konstanten Faktor) vollständig bestimmt wird.
Das Entwerfen von Filtern wie beispielsweise Tief-, Hoch-, Bandpassfilter ist direkt das Platzieren
\index{Tiefpassfilter}%
\index{Hochpassfilter}%
\index{Bandpassfilter}%
von Nullstellen und Polen in der komplexen Ebene,
genau in der Tradition von Eulers Produktansatz.
Die Stabilität eines Regelkreises hängt davon ab, ob alle Pole in der linken
Halbebene liegen, also eine Nullstellenbedingung im Sinne dieser Arbeit.
\index{Regelkreis}%
\index{Stabilitat@Stabilität}%

\subsubsection{Zahlentheorie und die riemannsche Vermutung}
Die riemannsche Zeta-Funktion $\zeta(s)$ besitzt eine Hadamard-Produktdarstellung
\index{Zeta-Funktion, riemannsch}%
\index{zeta(s)@$\zeta(s)$}%
über ihre sogenannten nichttrivialen Nullstellen.
Die Riemannsche Vermutung ist im Kern eine Aussage darüber, wo diese Nullstellen liegen.
Ihr Beweis oder ihre Widerlegung würde tiefe Aussagen über die Verteilung
der Primzahlen liefern.
Ohne die Theorie der ganzen Funktionen und ihren Produktsatz wäre das Problem
nicht einmal präzise formulierbar.
Für eine detaillierte Ausführung zu diesem Thema siehe Kapitel \ref{chapter:zeta}

\subsubsection{Quantenfeldtheorie und Regularisierung}
In der theoretischen Physik treten formal divergente unendliche Produkte auf,
etwa bei der Berechnung von Vakuumenergien.
\index{Vakuumenergie}%
Techniken wie die Zeta-Funktionen-Regularisierung und die dimensionale Regularisierung
\index{Regularisierung, dimensional}%
\index{dimensionale Regularisierung}%
lassen sich als kontrollierte Manipulationen des nullstellenfreien Faktors $e^{g(z)}$
im Weierstraß-Produkt verstehen.
Das bekannte Resultat $1 + 2 + 3 + \cdots = -1/12$ im Sinne der
Zeta-Regularisierung ist ein anschauliches, wenn auch vereinfachendes Beispiel
für diese Technik.

\subsubsection{Spezielle Funktionen in der Mathematik}
Die Gamma-Funktion selbst ist unverzichtbar in der Statistik
(Gamma-Verteilung, Beta-Funktion), in der Kombinatorik
(verallgemeinerte Binomialkoeffizienten) und in der Lösungstheorie
gewöhnlicher Differentialgleichungen (Bessel-Funktionen,
hypergeometrische Funktionen).
All diese Anwendungen bauen implizit auf ihrer Produktstruktur auf.

Der in dieser Arbeit entwickelte Gedankengang
von der elementaren Polynomfaktorisierung über Eulers kühnen heuristischen Schritt
bis zur rigorosen Weierstraß-Theorie 
illustriert exemplarisch, wie mathematische Intuition und strenge Analysis
zusammenwirken und wie abstrakte Theorie zu konkreten Anwendungen führt.
